%======================================================================
\section{Proofs for Section~\ref{sec:flat}: blocks and paths}\label{app:frames}
%======================================================================

Two further facts are used.
\begin{itemize}[leftmargin=2.8em]
\item[(F3)] \emph{Partial traces.} For a frame $F$ with $R=FF^*$ and operators $Y$ on the memory, $\Tr_{\rm mem}[(Y\otimes I)R(Y\otimes I)^*]=\Tr_{\rm mem}[(Y^*Y\otimes I)R]$;
the map $Y\mapsto\Tr_{\rm mem}[(Y\otimes I)R]$ is positive; and $\Tr_{\rm mem}[(P_p\otimes I)R]=\Tr_{H_p}R_p$.
\item[(F4)] \emph{Reduced states.} If $\rho(u)$ denotes the observer's reduced state of a vector $u$, then
$\frac12\nrm{\rho(u)-\rho(v)}_1\le\frac12\nrm{\ket u\bra u-\ket v\bra v}_1\le\frac12\nrm{u-v}(\nrm u+\nrm v)$, and isometries on systems the observer does not hold
do not change $\rho(\cdot)$. Errors are always compared in this form, so that vectors of different spaces, such as a physical device and a proof object, are
never subtracted.
\end{itemize}

\subsection{Concentration and export}

\begin{proof}[Proof of Lemma~\ref{lem:concentration}]
(a) Write $Z=\sum_pZ_p$ with $Z_p$ depending only on $V_p$. For each $p$ fix a Hermitian orthonormal basis $(T_{p,a})_a$ of $M_{D_p}$ for the Hilbert--Schmidt
inner product, so that $\sum_aT_{p,a}YT_{p,a}=\Tr(Y)I$ and $\sum_aT_{p,a}^2=D_pI$, and let $\partial_{p,a}$ be the derivative along $V_p\mapsto V_pe^{iuT_{p,a}}$ at
$u=0$. Haar measure is invariant under these flows, so $\E[(\partial f)h]=-\E[f\,\partial h]$ for smooth matrix-valued $f,h$. Put
$\mathsf L_p=\sum_a\partial_{p,a}^2$.

\emph{Generator.} With $X=V_p^*A_pV_p$, $\partial_{p,a}X=i[X,T_{p,a}]$ and $\partial_{p,a}^2X=-[T_{p,a},[T_{p,a},X]]$; summing, $\mathsf L_pX=-2D_pX+2\Tr(X)I$, and
$\Tr X=\Tr A_p$, so $\mathsf L_pZ_p=-2D_pZ_p$. Also $\partial_{p,a}Z=\partial_{p,a}Z_p$.

\emph{Carr\'e du champ.} $\partial_{p,a}Z=F_p^*(Y_{p,a}\otimes I)F_p$ with $Y_{p,a}=i(XT_{p,a}-T_{p,a}X)$ Hermitian, so
$\sum_a(\partial_{p,a}Z)^2=F_p^*[\sum_a(Y_{p,a}\otimes I)R_p(Y_{p,a}\otimes I)]F_p$. By $(P-P')R(P-P')^*\le2PRP^*+2P'RP'^*$,
$\sum_a(XT_a\otimes I)R_p(T_aX\otimes I)=(X\otimes I)(I\otimes\Tr_{H_p}R_p)(X\otimes I)\le\alpha_pX^2\otimes I$, and
$\sum_a(T_aX\otimes I)R_p(XT_a\otimes I)=I\otimes\Tr_{H_p}[(X\otimes I)R_p(X\otimes I)]\le\nrm X^2\alpha_pI$; the last inequality is tested on an exterior vector
$v$, for which $R_v=(I\otimes\bra v)R_p(I\otimes\ket v)\ge0$ satisfies $\Tr(XR_vX)\le\nrm X^2\Tr R_v$. Hence
$\sum_a(\partial_{p,a}Z)^2\le4\alpha_p\nrm{A_p}^2F_p^*F_p$, and this is $0$ when $A_p=0$.

\emph{Trace exponential.} Let $f(\theta)=\E\Tr e^{\theta Z}$. Integration by parts with $\mathsf L_pZ_p=-2D_pZ_p$ gives
$f'(\theta)=\sum_p\E\Tr(Z_pe^{\theta Z})=\sum_p(2D_p)^{-1}\sum_a\E\Tr[(\partial_{p,a}Z)(\partial_{p,a}e^{\theta Z})]$. With
$\partial e^{\theta Z}=\theta\int_0^1e^{u\theta Z}(\partial Z)e^{(1-u)\theta Z}du$ and $\int_0^1e^{u\lambda+(1-u)\mu}du\le(e^\lambda+e^\mu)/2$ in an eigenbasis of $Z$,
$\Tr[(\partial Z)(\partial e^{\theta Z})]\le\theta\Tr[(\partial Z)^2e^{\theta Z}]$. So
$f'(\theta)\le\theta\,\E\Tr[\sum_p(2\alpha_p\nrm{A_p}^2/D_p)F_p^*F_pe^{\theta Z}]\le2\theta\kappa g_0f(\theta)$, using
$\sum_{p:A_p\ne0}F_p^*F_p\le g_0I$ and $\Tr(Ye^{\theta Z})\le\nrm Y\Tr e^{\theta Z}$ for $Y\ge0$. Hence $f(\theta)\le Ne^{\theta^2\kappa g_0}$, and
$\Pr\{\lambda_{\max}(Z)\ge s\}\le e^{-\theta s}f(\theta)=Ne^{-s^2/(4\kappa g_0)}$ at $\theta=s/(2\kappa g_0)$; the same holds for $-Z$.

(b) Let $\widetilde T_{p,a}$ be $T_a$ on $H_p$ and $0$ elsewhere, and $\partial_{p,a}$ the derivative along $V_p\mapsto V_pe^{iuT_a}$; then
$\partial_{p,a}X=i[X,\widetilde T_{p,a}]$. Since $\sum_a\widetilde T_{p,a}Y\widetilde T_{p,a}=\Tr(P_pYP_p)P_p$ and $\sum_a\widetilde T_{p,a}^2=DP_p$,
$\sum_{p,a}[\widetilde T_{p,a},[\widetilde T_{p,a},X]]=2D(X-\sum_p\tau(X_{pp})P_p)$, and $\tau(X_{pp})=\tau(A_{pp})$ does not depend on $V$; so the Laplacian of
the product group satisfies $\mathsf LZ=-2DZ$. Here the equal dimensions are needed. With $Y_{p,a}=i(X\widetilde T_{p,a}-\widetilde T_{p,a}X)$,
\[
\sum_{p,a}(\partial_{p,a}Z)^2\le2F^*\Bigl[\sum_{p,a}(X\widetilde T_{p,a}\otimes I)R(\widetilde T_{p,a}X\otimes I)+\sum_{p,a}(\widetilde T_{p,a}X\otimes I)R(X\widetilde T_{p,a}\otimes I)\Bigr]F .
\]
The first sum is $(X\otimes I)(\sum_pP_p\otimes\Tr_{H_p}R_p)(X\otimes I)\le\alpha_*\nrm A^2I$. The second is $\sum_pP_p\otimes\Tr_{\rm mem}[(XP_pX\otimes I)R]$ by (F3),
and $XP_pX=V^*AP_pAV$ because $V$ commutes with $P_p$. The positive operator $AP_pA$ lives on the blocks of $N(p)=\{p':A_{p'p}\ne0\}$, so by Cauchy--Schwarz
$AP_pA\le|N(p)|\sum_{p'\in N(p)}P_{p'}AP_pAP_{p'}\le|N(p)|\sum_{p'\in N(p)}\nrm{A_{p'p}}^2P_{p'}=:\Lambda_p$, which commutes with $V$. Hence $V^*AP_pAV\le\Lambda_p$
and, by (F3), $\Tr_{\rm mem}[(XP_pX\otimes I)R]\le|N(p)|\sum_{p'}\nrm{A_{p'p}}^2\Tr_{H_{p'}}R_{p'}\le\Delta^2\nrm A^2\alpha_*I$. So
$\sum(\partial Z)^2\le2(1+\Delta^2)\nrm A^2\alpha_*F^*F$ for every $V$. As in (a), $f'(\theta)\le\theta(1+\Delta^2)\nrm A^2\alpha_*g\,f(\theta)/D$,
$f(\theta)\le N\exp(\theta^2(1+\Delta^2)\nrm A^2\alpha_*g/(2D))$, and the tail follows at $\theta=sD/((1+\Delta^2)\nrm A^2\alpha_*g)$.
\end{proof}

For a free permutation round with point blocks, the Hermitian and anti-Hermitian parts of $\lambda_g^*\lambda_h\otimes I$ ($g\ne h$) couple each block to at
most two blocks, so $\Delta=2$ there.

\begin{proof}[Proof of Lemma~\ref{lem:export}]
(a) Export and import rearrange the columns of each block isometrically: the port moves to the exterior, and a maximally entangled pair joins the new port
and a new exterior register. So the block Grams do not change. The new exterior marginal of block $p$ is $Y_p\otimes I_e/e$ with
$Y_p=\Tr_{\widetilde H_p}[(V'_p\otimes I)R_p(V'_p\otimes I)^*]$, a function of $V'_p$ alone, so $\alpha_p(F')=\nrm{Y_p}/e$, and $\E Y_p=(I_e/e)\otimes\Tr_{H_p}R_p$ has
norm $\alpha_p/e$. As in~\cite[Lemma~A.4]{DouglasMghirbiB}, the frame $\widehat F_p:\C^S\to H_p\otimes\C^N$ with entries
$\widehat F_{(q,y),\sigma}=\overline{F_{(q,\sigma),y}}$ satisfies $\Tr_{H_p}[(X\otimes I)R_p]=\widehat F_p^*(X^T\otimes I)\widehat F_p$, $\widehat F_p^*\widehat F_p=\Tr_{H_p}R_p$ and
$\alpha(\widehat F_p)=g_p$, and the block $(k,l)$ of $Y_p$ is $\Tr_{H_p}[(X_{kl}\otimes I)R_p]$ with $X_{kl}^T=\overline{V'_p}^{\,*}(I\otimes\ket k\bra l)\overline{V'_p}$, which is
Haar. Part (a) of Lemma~\ref{lem:concentration} for one block, applied to $\widehat F_p$ ($S$ columns) and to the Hermitian parts of $I\otimes\ket k\bra l$ at level
$z_p=2\sqrt{\alpha_pg_p\ln(8Se^2N)/D_p}$, fails for each of the $2e^2$ events with probability at most $2Se^{-\ln(8Se^2N)}=1/(4e^2N)$; so block $p$ fails with
probability at most $1/(2N)$, and some block with probability at most $\frac12$. Otherwise each block of $Y_p$ is within $2z_p$ of its mean, the $e\times e$ block
matrix within $2ez_p$, so $\nrm{Y_p}\le\alpha_p/e+2ez_p$ and $\alpha_p(F')\le\alpha_p/e^2+2z_p$. The second form is $2\sqrt{uv}\le\eta u+v/\eta$ with
$u=\alpha_p/e^2$ and $v=4e^2g_p\ln(8Se^2N)/D_p$.

(b) The block-$Q$ part of the frame after the round has $R'_Q=\sum_i(P'_QB_iV\otimes I)R(V^*B_i^*P'_Q\otimes I)$, so by (F3)
$\Tr_{H'_Q}R'_Q=\Tr_{\rm mem}[(V^*\Gamma_QV\otimes I)R]$. If $\Gamma_Q$ is block diagonal, $V^*\Gamma_QV=\bigoplus_pV_p^*(P_p\Gamma_QP_p)V_p\le\sum_p\nrm{P_p\Gamma_QP_p}P_p$,
and (F3) gives $\alpha_Q(F')\le\sum_p\nrm{P_p\Gamma_QP_p}\alpha_p(F)$.
\end{proof}

For a free permutation round with point blocks, a basis permutation before the round (absorbed in $F$) and a fixed partial injection $\iota$ of points after
it, $\Gamma_Q=\sum_gP_{(\iota\lambda_g)^{-1}(Q)}\otimes I_m$ is diagonal, $(\iota\lambda_g)^{-1}(Q)$ is at most one point for each $g$, and no two labels send a point
to the same point. So $\nrm{P_p\Gamma_QP_p}\le1$ for every point $p$, at most $r$ points have $\nrm{P_p\Gamma_QP_p}=1$, and $\alpha_Q(F')\le r\max_p\alpha_p(F)$: this is
the inflow factor of~\cite[Theorem~12.2]{DouglasMghirbiC}.

\subsection{Paths}

\begin{proof}[Proof of Lemma~\ref{lem:pathlaw}]
The matrix $M^{(b)}$ is the Gram matrix of the vectors $C^{(b)}_j$ in the normalized Hilbert--Schmidt space, so $M^{(b)}\ge0$, and
$\sum_bC^{(b)*}_iC^{(b)}_j=B_i^*(\sum_bJ_bJ_b^*)B_j=B_i^*B_j$ gives $\sum_bM^{(b)}=[\tau_k(B_i^*B_j)]=I_r$. With the Stinespring isometry $V:\xi\mapsto\sum_iA_i\xi\otimes\ket i$,
$\widetilde A^{(b)}:=\sum_{i,j}\bra iM^{(b)}\ket jA_i^*A_j=V^*(I\otimes M^{(b)})V$, so $\widetilde A^{(b)}\ge0$, $\sum_b\widetilde A^{(b)}=I$, and
$w(b)=\nrm{\widetilde A^{(b)}}\le\nrm{M^{(b)}}\le\Tr M^{(b)}$.

(a) $(\beta I-X_b)\otimes M^{(b)}$ and $(\beta I+X_b)\otimes M^{(b)}$ are positive; summing over $b$ and using $\sum_bI\otimes M^{(b)}\le I$ gives
$-\beta I\le\sum_bX_b\otimes M^{(b)}\le\beta I$.

(b) In the ideal experiment round $s$ applies $V_s:\xi\mapsto\sum_iA_i\xi\otimes\ket i_{E_s}$, and $E_s$ is never touched again. Let
$N_{\le s}=M^{(b_1)}\otimes\cdots\otimes M^{(b_s)}$ on $E_1,\dots,E_s$. Summing over the later letters (the matrices $M^{(b)}$ sum to $I$),
$\Pr(b_1,\dots,b_s)=\langle X|N_{\le s}|X\rangle$, and since nothing acts on $E_{\le s}$ after round $s$ this equals $\langle X_s|N_{\le s}|X_s\rangle$ for the global
vector $X_s$ right after round $s$. Write $X_s=(V_s\otimes I)Y_s$, with $Y_s$ the vector just before the use of round $s$. Then
$\Pr(b_{\le s})=\langle Y_s|\widetilde A^{(b_s)}\otimes N_{<s}|Y_s\rangle\le w(b_s)\langle Y_s|I\otimes N_{<s}|Y_s\rangle=w(b_s)\Pr(b_{<s})$, because
$(w(b_s)I-\widetilde A^{(b_s)})\otimes N_{<s}\ge0$ and $E_{<s}$ is untouched between round $s-1$ and $Y_s$. The same computation with $\sum_b\widetilde A^{(b)}=I$ shows that
$\Pr$ is a probability law. So the conditional law $q$ of $b_s$ given $b_{<s}$ satisfies $q_b\le w(b)$, and $\E[w(b_s)^{-\theta}\mid b_{<s}]\le\kappa_\theta$; by the tower
property $\E[W(p)^{-\theta}]\le\kappa_\theta^n$, and Markov's inequality gives $\Pr\{W(p)<2^{-\lambda}\}=\Pr\{W(p)^{-\theta}>2^{\theta\lambda}\}\le\kappa_\theta^n2^{-\theta\lambda}$.
Since $w\le1$, $W$ does not increase along extensions, so $P_t(\lambda)$ is prefix-closed; $W(p)$ depends only on the type of $p$; and
$\sum_{p\in[c]^t}W(p)=2^{tu'}$ with $W(p)\ge2^{-\lambda}$ on $P_t(\lambda)$ gives $|P_t(\lambda)|\le2^{\lambda+tu'}$.
\end{proof}

\begin{proof}[Proof of Theorem~\ref{thm:pathreg}]
\emph{The device.} Let $h=\log r$, $E_p=2^{\lceil h/2\rceil}$, $J_t=\lfloor th/2\rfloor$, $j_t=J_t-J_{t-1}$ and $e_t=2^{j_t}$ (the export schedule
of~\cite[Theorem~12.2]{DouglasMghirbiC}), $x=\log(32n/\epsilon^2)$, $\lambda=(n\log\kappa_\theta+x)/\theta$ and $P_t=P_t(\lambda)$. The memory consists of a register
with $2^{\lceil\lambda+nu'\rceil}\ge\max_t|P_t|$ values, holding $\mathrm{rank}_t(p)$, the rank of $p$ in a fixed enumeration of $P_t$; a register $\C^c$; $K$; and
$\C^m=\C^{m'}\otimes\C^{E_p}$, whose last factor is the port. So $Q=\lceil\lambda+nu'\rceil+\lceil\log c\rceil+\lceil\log k\rceil+\lceil\log m\rceil$. Initially the register
holds $\mathrm{rank}_0(\varnothing)$, $\C^c$ holds $\ket1$ and $K\otimes\C^m$ is maximally mixed. Round $t$ applies: (1) controlled on $\mathrm{rank}_{t-1}(p)$, a unitary
$V_{t,p}$ on $K\otimes\C^m$; (2) a unitary extension of $W\otimes I_m$ from $\C^d\otimes\ket1\otimes K\otimes\C^m$ to $\C^d\otimes\C^c\otimes K\otimes\C^m$, identifying
$J_bK$ with $\ket b\otimes K$, and the release of the output; (3) a permutation of the register values and $[c]$ that maps $(\mathrm{rank}_{t-1}(p),b)$ to
$(\mathrm{rank}_t(pb),1)$ whenever $pb\in P_t$ (an injective partial map, completed to a permutation); (4) controlled on $\mathrm{rank}_t(p')$, a unitary $V'_{t,p'}$; (5) if
$t<n$, the export of $j_t$ port qubits and the import of $j_t$ maximally mixed qubits. All unitaries are fixed in advance, so $C=0$ and $J=J_{n-1}$.

\emph{Proof object.} The untruncated device has memory $\bigoplus_{p\in[c]^t}H_p$, $H_p\cong K\otimes\C^m$ of dimension $D_b=km$, replaces step (3) by the relabelling
of block $p\otimes\ket b$ as block $pb$, and uses unitaries $V_{t,p}$, $V'_{t,p'}$ for all paths. It is an isometric evolution, and on the blocks of $P_{t-1}$ it
coincides with the device up to the relabelling of the register. Its frame $F_t:\C^{r^t}\to(\bigoplus_pH_p)\otimes\C^{S_t}$ has one column per Kraus history;
$F_0$ is one unit vector with $\alpha(F_0)=1/D_b$; and the column $(y,i)$ of block $pb$ after the round is $(C^{(b)}_i\otimes I_m\otimes I)(V_{t,p}\otimes I)F_{t-1,p}\ket y$.
Let $G_{t,\tau}=\sum_{p\text{ of type }\tau}F_{t,p}^*F_{t,p}$ and $\Gamma_{t,\tau}=\sum_{p\text{ of type }\tau}M^{(p_1)}\otimes\cdots\otimes M^{(p_t)}$.

\emph{Induction.} Put $\delta_1=\epsilon/2$, $\eta=\epsilon^2/(32n^2(n+1)^{c-1})$, $L_*=\ln(8D_br^nE_p^2c^n)$ and $A=4e(1+8E_p^2n^2L_*)$. For $0\le t\le n$ we show
(I1) $\nrm{F_t^*F_t-I}\le t\delta_1/n$; (I2) $\alpha_p(F_t)\le A/D_b$ for every $p$ (for $t\le n-1$); (I3) $\nrm{G_{t,\tau}-\Gamma_{t,\tau}}\le t\eta$ for every type
$\tau$. All hold at $t=0$. Assume them at $t-1$, so $g(F_{t-1})\le2$.

\emph{Choice of $V_t$.} Apply Lemma~\ref{lem:concentration}(a) to $F_{t-1}$ with the following operators; always $\nrm{A_p}\le r$, since $\nrm{B_i}^2$ and
$\nrm{C^{(b)}_i}^2$ are at most $\nrm{\sum_iB_i^*B_i}=r$, so $\kappa\le(A/D_b)r^2/D_b$ and $g_0\le2$. (E1) For each $(i,j)$, the Hermitian and anti-Hermitian parts of
$B_i^*B_j\otimes I_m$ in every block, at level $s_1=\delta_1/(2rn)$: the block $(i,j)$ of the full Gram matrix after the round is
$\sum_pF_p^*[(V_p^*(B_i^*B_j\otimes I_m)V_p)\otimes I]F_p$, because children of different letters or different parents are orthogonal, and its mean is
$\tau(B_i^*B_j)F^*F=\delta_{ij}F^*F$; on success the $r\times r$ block matrix deviates by at most $2rs_1=\delta_1/n$, which gives (I1). (E2) For each type $\tau'$ of
length $t$, each letter $b$, each $(i,j)$ and the Hermitian or anti-Hermitian part of $C^{(b)*}_iC^{(b)}_j\otimes I_m$ on the blocks of type $\tau'-e_b$ ($A_p=0$ elsewhere),
at level $s_3=\eta/(2cr)$. The new type Gram is $G'_{\tau'}=\sum_b\sum_{p\in\tau'-e_b}F_p^*[(V_p^*X^{(b)}V_p)\otimes I]F_p$ with
$X^{(b)}=[C^{(b)*}_iC^{(b)}_j\otimes I_m]_{ij}$, whose mean is $\sum_bG_{t-1,\tau'-e_b}\otimes M^{(b)}$ by Schur's lemma; on success it deviates from that mean by at most
$2crs_3=\eta$, and Lemma~\ref{lem:pathlaw}(a) with $X_b=G_{t-1,\tau'-e_b}-\Gamma_{t-1,\tau'-e_b}$ gives $\nrm{G_{t,\tau'}-\Gamma_{t,\tau'}}\le(t-1)\eta+\eta$, which is (I3).
Each event fails with probability at most $2r^n\exp[-s^2D_b^2/(8Ar^2)]$; there are $N_E=2r^2(1+c(n+1)^{c-1})$ events and $s_3\le s_1$, so all succeed with
probability at least $\frac12$ if
\begin{equation}\label{eq:Db}
D_b^2\ \ge\ 8Ar^2\ln(4r^nN_E)/s_3^2,\qquad s_3=\epsilon^2/(64crn^2(n+1)^{c-1}).
\end{equation}
Fix such a $V_t$.

\emph{Inflow.} By (F3), for the child block $pb$, $\Tr_{H_{pb}}(F'_{pb}F'^*_{pb})=\Tr_{H_p}[((V_p^*\Gamma_bV_p)\otimes I)R_p]$ with
$\Gamma_b=\sum_iC^{(b)*}_iC^{(b)}_i\otimes I_m\le\sum_iB_i^*B_i\otimes I_m=rI$, so $\alpha_{pb}\le r\,a_{t-1}$ with $a_{t-1}:=\max_p\alpha_p(F_{t-1})$. Each child has a single
parent; this is where the tree structure of paths is used.

\emph{Choice of $V'_t$.} Lemma~\ref{lem:export}(a) with $e=e_t$, $\eta=1/n$, $g_p\le2$, at most $c^t$ blocks and $S_te_t^2\le D_br^nE_p^2$ gives
$a_t\le(1+1/n)r\,a_{t-1}/e_t^2+8E_p^2nL_*/D_b$; in rounds with $j_t=0$, $a_t\le r\,a_{t-1}$ directly. With $c_t=2^{2J_t}r^{-t}\in(\frac14,1]$ and $c_tr/e_t^2=c_{t-1}$,
$c_ta_t\le(1+1/n)c_{t-1}a_{t-1}+8E_p^2nL_*/D_b$, so $c_ta_t\le(1+1/n)^t(1+8E_p^2n^2L_*)/D_b$ and $a_t\le A/D_b$, which is (I2).

\emph{Error.} Purify the observer. The device acts on the input and the memory only through $W$, so the untruncated final vector is
$\tilde x_n=\sum_y\psi_y\otimes F_n\ket y=(I\otimes F_n)X$, with the observer's vectors $\psi_y$ of the ideal experiment and $X=\sum_y\psi_y\otimes\ket y$ a unit vector, and
$\nrm{\tilde x_n}=1$. Let $\Pi$ project onto the blocks of $P_n$, $\omega=\nrm{(I-\Pi)\tilde x_n}^2$, $x_n$ the physical final vector and $T$ the polar isometry of $F_n$.
\begin{enumerate}[label=\arabic*.,itemsep=1pt]
\item \emph{Polar step.} $\nrm{\tilde x_n-(I\otimes T)X}\le\nrm{(F_n^*F_n)^{1/2}-I}\le\nrm{F_n^*F_n-I}\le\delta_1$, since $|\sqrt u-1|\le|u-1|$ for $u\ge0$.
\item \emph{Truncation.} The untruncated dynamics never changes the prefix of a path, the observer does not act on the memory, and $P_t$ is prefix-closed; so
the projection onto ``first exit from $P$ at time $t$'' commutes with all later rounds, and with $\Pi_{{\rm fe}=t}$ that projection,
$\sum_t\nrm{\Pi_{{\rm fe}=t}\tilde x_n}^2=\omega$. Let $g_t$ be the good part of $\tilde x_t$ (blocks of $P_t$), embedded in the physical memory, and $\mathcal U_t$ the
physical round together with the observer's operation. On the good part the physical round is the untruncated one, and the permutation of step (3) maps good
inputs bijectively onto good outputs, so $\mathcal U_tg_{t-1}=g_t+d_t$ with $\nrm{d_t}=\nrm{\Pi_{{\rm fe}=t}\tilde x_t}=\nrm{\Pi_{{\rm fe}=t}\tilde x_n}$. From
$x_t-g_t=\mathcal U_t(x_{t-1}-g_{t-1})+d_t$ and $x_0=g_0$, $\nrm{x_n-g_n}\le\sum_t\nrm{d_t}\le\sqrt{n\omega}$. Discarded branches may re-enter good register values;
the bound does not care.
\item \emph{Discarded weight.} $P_n$ is a union of types, so $\omega=\sum_{\tau\notin P_n}\langle X|I\otimes G_{n,\tau}|X\rangle\le\sum_{\tau\notin P_n}\langle X|I\otimes\Gamma_{n,\tau}|X\rangle+|\mathcal T_n|n\eta$,
using $|\langle X|I\otimes Y|X\rangle|\le\nrm Y$ and (I3); here $|\mathcal T_n|\le(n+1)^{c-1}$ is the number of types, so $|\mathcal T_n|n\eta\le\epsilon^2/(32n)$. The first
term is the probability under the path law that $W(p)<2^{-\lambda}$, at most $\kappa_\theta^n2^{-\theta\lambda}=2^{-x}=\epsilon^2/(32n)$ by Lemma~\ref{lem:pathlaw}(b). So
$\omega\le\epsilon^2/(16n)$.
\item By (F4), with $\nrm{x_n}=\nrm{\tilde x_n}=1$ and $\nrm{g_n}\le1$, the observer's state from $x_n$ is within $\sqrt{n\omega}$ of that from $g_n$, which equals that from
$\Pi\tilde x_n$; that is within $\sqrt\omega$ of the state from $\tilde x_n$, which is within $\delta_1$ of the ideal state from $(I\otimes T)X$. The error is at most
$\delta_1+(1+\sqrt n)\sqrt\omega\le\epsilon/2+(1+\sqrt n)\epsilon/(4\sqrt n)<\epsilon$.
\end{enumerate}

\emph{Memory.} $Q\le\lambda+nu'+\log c+\log k+\log m+4$. Condition~\eqref{eq:Db} reads $D_b\ge K_{r,c}n^{c+3}\epsilon^{-2}(\ln D_b+n)^{1/2}$ up to constants, since
$A=O(n^2L_*)$ and $L_*=O_r(n+\ln D_b)$; take $m'$ least with $D_b=km'E_p$ satisfying it. Then $\log m=\log(D_b/k)=O_r(\log(n/\epsilon)+\log(2+\log k))$, as in the
proof of~\cite[Theorem~12.2]{DouglasMghirbiC}.
\end{proof}

\subsection{Flat rounds}

\begin{proof}[Proof of Lemma~\ref{lem:top}]
Each $B_gB_g^*/k\le I/k$, so the eigenvalues $\lambda_1\ge\lambda_2\ge\cdots$ of $\bar\tau$ are at most $1/k$, and $a=\sum_i\lambda_i\log(1/(k\lambda_i))$ with nonnegative
terms. There are at least $k$ positive eigenvalues, so $x:=\lambda_k>0$; put $y=kx\in(0,1]$. For $i>k$, $\lambda_i\le x$, so the tail contributes at least $w\log(1/y)$.
For $i\le k$, $\log(1/(k\lambda_i))\ge(1-k\lambda_i)/\ln2$ and $\lambda_i\ge x$, so the top contributes at least $(x/\ln2)\sum_{i\le k}(1-k\lambda_i)=(x/\ln2)(k-k(1-w))=yw/\ln2$.
Hence $a\ge(w/\ln2)(y-\ln y)\ge w/\ln2$.
\end{proof}

\begin{proof}[Proof of Theorem~\ref{thm:flat-zero}]
(a) By Lemma~\ref{lem:zero-leak} with $u=0$, the round restricted to $K=\supp\rho$ is an exact rectangular dilation $(B_g)$, and since the $A_g$ have disjoint supports,
$T(\rho)=\sum_g\mu_gB_gB_g^*/k=:\bar\tau$, so $a=S(\bar\tau)-\log k$. Restrict $\C^{k'}$ to the span of the ranges ($k'\le rk$) and choose the pieces: $J_1K$ the
top-$k$ eigenspace of $\bar\tau$, and the complement, padded, cut into $k$-dimensional pieces, $c\le r+1$ in all. Then
$\bra gM^{(1)}\ket g=\Tr(P_{J_1K}\omega_g)$ with $\omega_g=B_gB_g^*/k$, and
\[
\beta:=\sum_{b\ge2}w(b)\le\sum_{b\ge2}\Tr M^{(b)}=\sum_g(1-\bra gM^{(1)}\ket g)\le\mu_{\min}^{-1}\sum_g\mu_g(1-\bra gM^{(1)}\ket g)\le d\,a\ln2
\]
by Lemma~\ref{lem:top}, since $\mu_{\min}\ge1/d$. As $w(1)\le1$, $u'\le\log(1+\beta)\le1.443\beta$. For $q\le w$, $\sum_bq_bw(b)^{-\theta}\le\sum_bw(b)^{1-\theta}\le1+(c-1)^\theta\beta^{1-\theta}$
by H\"older's inequality over $b\ge2$, so $\kappa_\theta\le1+(c-1)^\theta\beta^{1-\theta}$. Take $\theta=1/\log(1/\bar\beta)$ with $\bar\beta=\max(\beta,\epsilon^2/n^2)\le\frac14$;
this needs $a\le1/(4d\ln2)$. Then $\theta\le\frac12$ and $\bar\beta^{-\theta}=2$, so $(c-1)^\theta\beta^{1-\theta}\le2\sqrt{c-1}\,\bar\beta$ and
$\log\kappa_\theta/\theta\le2.885\sqrt{c-1}\,\bar\beta\log(1/\bar\beta)$. Also $\log(32n/\epsilon^2)/\theta=\log(1/\bar\beta)\log(32n/\epsilon^2)\le2\log(n/\epsilon)\log(32n/\epsilon^2)$,
and if $\bar\beta=\epsilon^2/n^2$ then $n\bar\beta\log(1/\bar\beta)=2(\epsilon^2/n)\log(n/\epsilon)\le2$. Theorem~\ref{thm:pathreg} gives the first bound, since $x\log(1/x)$
increases on $(0,1/e)$ and $\beta\le d\,a\ln2$. For the second: if $a\ge\epsilon^2/n^2$ then $\log(e/a)\le1.443+2\log(n/\epsilon)$, and if $a<\epsilon^2/n^2$ then
$n\,a\log(e/a)\le(\epsilon^2/n)(1.443+2\log(n/\epsilon))\le2.45$; so $n\,a\log(e/a)\le2.45+n\,a(1.443+2\log(n/\epsilon))$, and every term is at most
$K\log^2(n/\epsilon)(1+\log k+n\,a)$. If $a>1/(4d\ln2)$, Theorem~\ref{thm:robust}(a) with $G=I$, which is~\cite[Theorem~12.2]{DouglasMghirbiC}, gives $Q\le\log k+n\log(k'/k)+O_r(\log(n/\epsilon)+\log(2+\log k))$ with
$\log(k'/k)\le\log r\le4d\ln2\log r\cdot a$.

(b) The changes to the proof of Theorem~\ref{thm:pathreg} are those of Theorem~\ref{thm:robust}(a). (I1) becomes $\nrm{F_t^*F_t-G^{\otimes t}}\le e^{1/2}t\delta_1/n$:
the mean of the block $(i,j)$ after a round is $\tau(B_i^*B_j)F^*F$, so the deviation $e_t$ from $G^{\otimes t}$ satisfies $e_t\le\nrm Ge_{t-1}+\delta_1/n$, and
$\nrm G^n\le e^{n\gamma}\le e^{1/2}$; then $g(F_t)\le4$, and $2$ is replaced by $4$ in~\eqref{eq:Db} and in Lemma~\ref{lem:export}(a). (I2) is unchanged, since
$\sum_iB_i^*B_i=rI$ by Lemma~\ref{lem:dilation}(b). (I3) becomes $\nrm{G_{t,\tau}-\Gamma_{t,\tau}}\le e^{1/2}t\eta$; we run the construction with $\eta/\sqrt e$ in place of $\eta$, which
restores the discarded-weight bound $\omega\le\epsilon^2/(16n)$ and changes only constants. This uses Lemma~\ref{lem:pathlaw}(a) with
$\sum_bM^{(b)}=G$, which gives $\nrm{\sum_bX_b\otimes M^{(b)}}\le\beta\nrm G$. Applying $\id\otimes\tau_K$ to $W^*W=I$ gives
$\sum_b\widetilde A^{(b)}=\sum_{i,j}G_{ij}A_i^*A_j=I$, so the path law is still a probability law and Lemma~\ref{lem:pathlaw}(b) holds verbatim. In the polar step,
$\Gamma=G^{\otimes n}\ge(1-\gamma)^nI\ge e^{-1}I$, and as in Theorem~\ref{thm:robust}(a), $\nrm{(F_n^*F_n)^{1/2}-\Gamma^{1/2}}\le e^{1/2}\delta_1/e^{-1/2}=e\delta_1$; so
$\tilde x_n$ is within $e\delta_1$ of $(I\otimes T)(I\otimes\Gamma^{1/2})X$, the final vector of the experiment that applies $\Phi_W(X)=\sum G_{ji}A_iXA_j^*$ in every round, and a
hybrid over the rounds costs $n\,\ddia(\Phi_W,\Phi)\le\frac n2\gamma$. With $\delta_1=\epsilon/(2e)$ the error is at most $\epsilon+\frac n2\gamma$.

(c) By Lemma~\ref{lem:zero-leak} the round gives isometries $B_g=X_g|_K$ forming a gadget family with flat Gram $G$, $G_{gg}=1$ and $|G_{gh}|\le c$. By
Theorem~\ref{thm:robust}(b) the family $(B_g^{\otimes t})$ is a gadget family with Gram entries $G_{gh}^t$, so $\gamma_t:=\nrm{G^{(t)}-I}\le(r-1)c^t\le\epsilon/(4n)$,
$n\gamma_t\le\frac12$ and $\frac n2\gamma_t\le\epsilon/8$. Its flat entropy production $a_t=S(\sum_g\mu_g\omega_g^{\otimes t})-t\log k$ is at most $t\,a$ by subadditivity,
since every one-factor marginal of $\sum_g\mu_g\omega_g^{\otimes t}$ is $\bar\tau$. Restrict the codomain to the span of the ranges of the $B_g^{\otimes t}$, of dimension at
most $rk^t$, and take the pieces of (a) for the $t$-fold family; Lemma~\ref{lem:top} applies to it, so $\beta_t\le d\,a_t\ln2\le d\,t\,a\ln2$. Part (b) with the choice of
$\theta$ of (a) gives the second bound when $d\,t\,a\ln2\le\frac14$; otherwise Theorem~\ref{thm:robust}(c), run at error $2\epsilon/(e^2+1)$ so that
its error is at most $2\epsilon$, gives $Q\le t'\log k+nt'\log(k'/k)+O_r(\cdot)$ with its own $t'\le t+2$ and
$\log(k'/k)\le\log r\le4d\,t\ln2\log r\cdot a$. For the polylogarithmic form, $t\le\log(4rn/\epsilon)+1\le3\log(n/\epsilon)+\log(8r)$ for $n\ge2$ and $\epsilon<1$, and
$ta\log(e/(ta))\le t\,a\log(e/a)$.
\end{proof}
